\section{Experimental setup}\label{sec:setup}
Phase~1 is development: the method was changed and judged openly on a small set of instances, so every number in
Section~\ref{sec:results} is descriptive. Table~\ref{tab:terms} defines the terms used from here on.

\begin{table}[t]\centering\footnotesize
\caption{Terms used in this report.}\label{tab:terms}
\begin{tabular}{>{\raggedright}p{2.6cm}>{\raggedright\arraybackslash}p{13.0cm}}
\toprule
term & meaning \\
\midrule
primal integral $P$ & $P(120)$: the gap $|v-r|/\max(|v|,|r|)$ of the best checked solution value $v$ to the reference
  $r$ (1 before the first solution or if the signs differ), averaged over the 120-s budget. \\
coupled system & the method of Section~\ref{sec:method}; its $P$ is that of the best solution held by any of its
  processes. \\
cuOpt alone & cuOpt~26.08 on the coupled system's share of the machine (5 of its 20 CPU cores and the GPU), run in
  the same window. \\
SCIP alone & SCIP~10 on one CPU core, two seeds, run on 6--7 October. \\
bounded difference $d$ & per instance, on seed means: $(P_\text{alone} - P_\text{coupled}) / \max(P_\text{alone},
  P_\text{coupled}, 0.001)$, in $[-1, 1]$; positive when the coupled system has the lower $P$. \\
hard set & the instances where a solver alone has $P \ge 0.5$: 11 for cuOpt, fixed on earlier runs of cuOpt alone
  with three cores; 15 for SCIP, on the compared runs (Section~\ref{sec:limits}). \\
hint & mean $P$ of the coupled system over that of cuOpt alone, on cuOpt's hard set; present if at most 0.95. It only
  decides whether testing continues. \\
guard & mean $d$ against cuOpt alone over all 17 instances; holds if at least $-0.05$. \\
harm check & per instance, the final configuration against the one before LP decoding and the yielding wait (13
  October); flags if their $d_i$ (as $d$) is below minus the instance's seed-to-seed band (at least 0.05). \\
\bottomrule
\end{tabular}
\end{table}

\paragraph{Machine.} One NVIDIA DGX Spark: a GB10 with 10 Cortex-X925 and 10 Cortex-A725 CPU cores, a Blackwell GPU and
121\,GiB of memory shared by CPU and GPU. The software is SCIP~10 through PySCIPOpt 6.2.1, HiGHS 1.15.1, cuOpt~26.08
and CuPy 14.2.0 on CUDA~13. No other GPU process ran beside the coupled system. Runs start only below temperature
gates, and a run whose GPU clock drops is voided and repeated (Appendix~\ref{app:harness}).

\paragraph{Instances.} We used 17 public instances (Appendix~\ref{app:c1}). Nine were chosen after a screening run of
SCIP alone. Seven are MIPLIB 2017 instances~\citep{gleixner2021miplib} on which SCIP held a large gap or found no
solution; one is a PGLib-UC unit-commitment case~\citep{pglibuc2019,knueven2020uc} and one an ML4CO item-placement
instance~\citep{gasse2022ml4co}. Eight MIPLIB instances from other families were added later. All 17 were chosen by
looking at results, the method was tuned on them, and they are excluded from later test sets. References are MIPLIB's
optimal values and provisional best known values for the other two.

\paragraph{Arms.} Several processes can do better than one simply by being several. cuOpt alone therefore gets the
coupled system's whole machine share: two Cortex-X925 cores, three Cortex-A725 cores and the GPU. Both arms of each
instance and seed ran back to back in one window, under the same memory caps. SCIP alone had fewer resources: one core,
no GPU, run on 6--7 October. Only solutions that pass an independent exact checker count in $P$.

\paragraph{Statistics.} We report mean $d$ with a 95\,\% percentile bootstrap interval over instances (10{,}000
resamples)~\citep{dicicco1996bootstrap} and the Wilcoxon signed-rank test on $d$ (ties dropped). The checkpoint
result uses seeds 0--3 (SCIP alone: 0--1); development tests used one to six seeds.

\paragraph{How changes were judged.} Each change was tested against the unchanged configuration in the same window,
under a dated record with a written reading rule. It was kept only if its checks all passed: validity, the intended
mechanism, no harm per instance, and the hint. Most rules were fixed before any run was read; Appendix~\ref{app:devrecord} lists the
exceptions. A crashed run counts for its arm, scored up to the crash.

\paragraph{The checkpoint.} At the end of phase~1 the author decides whether to go on to unseen instances, reading the
hint, guard and harm check under rules fixed before any checkpoint run was read.
\FloatBarrier
